\section{The two quartic Tate residues for the actual conductor-sixteen characters}
\label{sec:conductor16-quartic-residue-strict}

\paragraph{Scope and class input.}
This proof uses AC-G147-class only through
$\operatorname{Cl}(L_1)[269]=\operatorname{Cl}(L_2)[269]=0$
for its two explicitly identified real-quartic composita.
The analytic selectors are not class-number assumptions.
The local unit calculations below are separate exact computations,
with every rational denominator retained.
No automorphic point, auxiliary host, model, or mixed realization is used.
The coefficient ring remains the original ring throughout.
The single prime of ramification degree twelve below replaces the
two-prime, ramification-degree-six situation of AC-G144.

Put $p=269$, $K=\mathbb Q(\omega)$, $t=\sqrt[3]{p}$ with its
real embedding, and let $\chi:G_K\to\mu_3$ describe $K(t)/K$.
Use the arithmetic conductor-sixteen character
\[
 \ker\phi=\{1,7\},\qquad \phi(13)=i,\qquad
 \phi(5)=-i,\qquad i\bmod p=82 .
 \tag{C16Q.1}
\]
Thus $\phi$ is odd and $\phi^2=\chi_8$.
Let $c_p$ be the actual cyclotomic character, let $\omega_p$ be
its Teichmuller part, and set
\[
                  \psi_p=\omega_p^{67},\qquad
                  \xi_1=\phi\psi_p,\quad \xi_2=\phi\psi_p^{-1}.
 \tag{C16Q.2}
\]
Both $\psi_p$ and $\phi$ are odd; $\psi_p(2)=-i$.
The two characters $\xi_1,\xi_2$ are even quartic characters.
They are used only in Kummer splitting fields, not as additional
finite labels in the original host.

For any discrete coefficient $D$ write
\[
 \begin{gathered}
 \operatorname{Sel}_{\rm str}(K,D)=
 \ker\bigl(H^1_{\rm cts}(G_K,D)\longrightarrow\mathcal L(D)\bigr),\\
 \mathcal L(D)=
 \prod_{v\mid p}H^1(G_{K_v},D)\ \times\
 \prod_{v\nmid p,\infty}H^1(I_v,D).
 \end{gathered}
 \tag{C16Q.3}
\]
In particular the condition at $p$ is zero in full local cohomology.

\begin{theorem}[AC-G147: the two quartic Tate residues on the original coefficient ring]
\label{AC-G147}
Let $E/\mathbb Q_p$ be finite, let $O$ be its original ring of integers,
and suppose that $O$ contains the values of $\chi$.
For
\[
 \begin{gathered}
 j\in\{1,2\},\quad \sigma\in\{1,-1\},\quad m\in\mathbb Z,\quad e\geq1,\\
 \kappa:G_K\longrightarrow1+\mathfrak m_O
       \quad\text{continuous and unramified away from }p,
 \end{gathered}
 \tag{C16Q.4}
\]
put
\[
       D=(E/O)(\chi^{-j}\phi^{-\sigma}
                    \omega_p^{-m}\kappa^{-1}c_p^{-e}).
 \tag{C16Q.5}
\]
If $e+m\equiv66$ or $200\bmod268$, then
\[
                       \operatorname{Sel}_{\rm str}(K,D)=0.
 \tag{C16Q.6}
\]
The residual strict kernel also vanishes over every finite field
containing its character values.
There is no absolute ramification restriction on $O$.
For $m=0,\kappa=1$, the four finite labels are exactly
$\chi^{-j}\phi^{-\sigma}$ of AC-K25, at the two new Tate residues.
Neither $\psi_p$ nor its quadratic square is inserted into the host.
\end{theorem}

\begin{proof}
\proofstep{1. Identify the two character fields, including their orientation.}
The exact cyclotomic construction in AC-G147-class gives
\[
 \begin{gathered}
 F_1=\mathbb Q(w_1),\qquad w_1^2=538+23\sqrt{538},\\
 w_1^4-1076w_1^2+4842=0,\\
 F_2=\mathbb Q(w_2),\qquad w_2^2=538+3\sqrt{538},\\
 w_2^4-1076w_2^2+284602=0.
 \end{gathered}
 \tag{C16Q.7}
\]
They are respectively the fields cut out by $\xi_1$ and $\xi_2$.
We recall the exact identification rather than infer it from
equal discriminants.
In $\mathbb Q(\zeta_p)$, let $\eta_r$ be the sum over the quartic
residue class $2^r(\mathbb F_p^\times)^4$, and put
$v=\eta_0-\eta_2$, $v'=\eta_1-\eta_3$.
With the oriented quadratic Gauss sum $s=\sqrt p$, the retained
finite identities are
\[
             2v^2=-p-13s,\qquad vv'=-5s,\qquad s^2=p.
 \tag{C16Q.8}
\]
Let $w_0=\zeta_{16}^{3}-\zeta_{16}^{-3}$ and let $b_0$ act by
$\zeta_{16}\mapsto\zeta_{16}^{13}$, so $\phi(b_0)=i$.
The two elements $w_0v+b_0(w_0)v'$ and
$w_0v-b_0(w_0)v'$ have the squares in (C16Q.7).
The simultaneous actions $(b_0,\zeta_p\mapsto\zeta_p^2)$ and
$(b_0,\zeta_p\mapsto\zeta_p^{2^{-1}})$ fix the respective elements.
Their character values under $\phi\psi_p$ and $\phi\psi_p^{-1}$
are one.  These kernels identify $F_1,F_2$ with the claimed fields.
They are distinct cyclic quartic fields, not two generators of one field.
Both have the unique quadratic subfield $\mathbb Q(\sqrt{538})$.

For each index $i$ put
\[
                    L_i=F_i(t),\qquad M_i=K L_i.
 \tag{C16Q.9}
\]
The fields $F_i$ are unramified at three.
The only nontrivial abelian subfield of the $S_3$ splitting field
$K(t)$ is $K$, which ramifies at three.
Therefore $F_i\cap K(t)=\mathbb Q$, and
\[
 \operatorname{Gal}(M_i/\mathbb Q)=S_3\times C_4,\qquad
 \Delta_i:=\operatorname{Gal}(M_i/K)=C_3\times C_4.
 \tag{C16Q.10}
\]
In particular $[L_i:\mathbb Q]=12$, $[M_i:\mathbb Q]=24$,
and $\mu_p\not\subset M_i$, since $268$ does not divide $24$.
Choose the generator $b$ induced by $b_0$ and identity on the
$p$-cyclotomic factor.  It fixes $K,t$, has $\xi_i(b)=i$, and
the exact global formulas are
\[
 \begin{array}{c|c|c}
 i&c_i& b(w_i)\\ \hline
 1&18&(w_1^3-1067w_1)/69\\
 2&1058&(w_2^3-547w_2)/69 .
 \end{array}
 \tag{C16Q.11}
\]
Here both defining polynomials are
$w_i^4-4p w_i^2+c_i p=0$.
The identities $b^2(w_i)=-w_i$ follow in these exact quotient rings.

\proofstep{2. A single prime and an exact uniformizer model.}
Fix $i$, abbreviate $c=c_i,w=w_i$, and let $k_i$ be $1067$ or $547$
as in (C16Q.11).
The quartic polynomial is Eisenstein at $p$, so $F_i$ is totally
ramified of degree four there.
The extension generated by $t$ is totally ramified of degree three.
Their compositum has ramification degree divisible by both three
and four, and degree at most twelve.
It therefore has degree and ramification degree twelve, with
residue field $\mathbb F_p$.
The field $K_p$ is unramified quadratic.
Consequently $M_i$ has exactly one prime above $p$, of local data
\[
                         e_p=12,\qquad f_p=2.
 \tag{C16Q.12}
\]
Indeed its completion already has degree twenty-four, the full
global degree.  In particular the whole $\Delta_i$ acts on this
single completion; no permutation of different completions is suppressed.

Set
\[
                         u=\frac tw.
 \tag{C16Q.13}
\]
Its valuation is one when $v_u(p)=12$, since $v_u(t)=4$ and
$v_u(w)=3$.
Eliminating $t,w$ gives the Eisenstein equation
\[
                  c^3u^{12}+p(12c-64p)u^6+p=0.
 \tag{C16Q.14}
\]
Conversely the exact identities
\[
 \begin{gathered}
 B=16p-c,\qquad Q=c^2-48cp+256p^2,\\
 w=\frac{Q u^3-4c^3u^9}{B},\qquad t=u w
 \end{gathered}
 \tag{C16Q.15}
\]
hold in its rational quotient ring.
Here $B$ and $c$ are $p$-adic units.
They recover $t^3=p$ and $w^4-4p w^2+cp=0$, so this is
the same field used in the class theorem, not a numerical analogue.
Equivalently, with $z=w/u^3$, the identities
$z=4u^6z^2-c$ and $p=u^{12}z^3$ imply both displayed formulas.

The ring of integers of the completion of $L_i$ is
$\mathbb Z_p[u]$; after adjoining $K_p$ it is
$\mathcal O_{K_p}[u]$.
For this assertion one may use the unique expansion in the basis
$1,u,\ldots,u^{11}$: the valuations of its nonzero coefficient
terms are distinct modulo twelve, so an element is integral
exactly when all rational $p$-adic coefficients are integral.
This also explains the integrality test in the exact verifier.

Modulo $u^{12}$, $p$ vanishes, and (C16Q.15) gives
\[
             w=-c u^3+4c^2u^9,\qquad
             t=-c u^4+4c^2u^{10}.
 \tag{C16Q.16}
\]
Since $b(t)=t$, formula (C16Q.11) then gives
\[
 \begin{gathered}
 b(u)=a_i u+a_i c_i^2 k_i^{-1}u^7\pmod{u^{12}},
             \qquad a_i=-69k_i^{-1}\in\mathbb F_p,\\
 b(u)=187u+262u^7\quad(i=1),\qquad
 b(u)=82u+262u^7\quad(i=2).
 \end{gathered}
 \tag{C16Q.17}
\]
Direct substitution gives $b^2(u)=-u$ at this precision.
If $d$ generates the cubic factor, then $d(u)=\omega u$.
These are commuting local actions on the same completion.

\proofstep{3. Residual Kummer descent with the inverse character.}
Let $k=O/\mathfrak m_O$.
Fourth roots of unity already belong to $\mathbb Q_p$, so
no coefficient extension for $\phi$ is required.
Reduction of the character of (C16Q.5) gives
\[
       a=\bar\chi^{-j}\bar\phi^{-\sigma}\bar c_p^{-(e+m)},
       \qquad \eta=a/\bar c_p .
 \tag{C16Q.18}
\]
The exact quartic factor of $\eta$ is as follows:
\[
 \begin{array}{c|c|c|c}
 e+m\bmod268&\sigma&\text{quartic factor of }\eta&\text{field}\\ \hline
 66&+1&\bar\xi_1^{-1}&F_1\\
 66&-1&\bar\xi_2&F_2\\
 200&+1&\bar\xi_2^{-1}&F_2\\
 200&-1&\bar\xi_1&F_1 .
 \end{array}
 \tag{C16Q.19}
\]
For example $-(66+1)=-67$ and $-(200+1)=67\bmod268$.
The cubic factor is $\bar\chi^{-j}$ in every row.
Thus $\eta$ is faithful on the corresponding $\Delta_i$, and
its splitting field is $M_i$.

Prime-to-$p$ inflation--restriction, globally and on every local
subgroup, identifies the residual strict kernel with the
$\eta^{-1}$ part of the full Kummer localization kernel on $M_i$.
To make both the orientation and the local condition explicit, put
\[
 \begin{gathered}
 V_{M,S}=\{x\in M^\times:p\mid v(x)\text{ for }v\nmid p\}/M^{\times p},\\
 V_M=\{x\in M^\times:p\mid v(x)\text{ at every finite }v\}/M^{\times p},\\
 \mathcal P_M=M_p^\times/M_p^{\times p}.
 \end{gathered}
 \tag{C16Q.20}
\]
Then the kernel is
\[
 e_{\eta^{-1}}\ker\left(
       V_{M,S}\otimes_{\mathbb F_p}k
        \longrightarrow\mathcal P_M\otimes_{\mathbb F_p}k\right).
 \tag{C16Q.21}
\]
Over $M$, the coefficient is $k(1)$, and the extra
$\Delta_i$-twist on its Kummer cohomology is $\eta$.
Taking invariants therefore selects $\eta^{-1}$.
Away from $p$, inertia zero is exactly divisibility of the
Kummer valuation by $p$.
A local $p$th power at $p$ has its valuation divisible by $p$ too.
The kernel in (C16Q.21) consequently belongs to $V_M$.
This uses full local powers, not a valuation-only condition.

The all-primes valuation sequence is
\[
 0\longrightarrow\mathcal O_M^\times/\mathcal O_M^{\times p}
 \longrightarrow V_M\longrightarrow\operatorname{Cl}(M)[p]
 \longrightarrow0.
 \tag{C16Q.22}
\]
Only $\operatorname{Cl}(L_i)[p]=0$ is needed from AC-G147-class.
Let $s$ be the reflection in the $S_3$ factor fixing $L_i$.
On $\operatorname{Cl}(M_i)[p]$, norm followed by extension yields
$1+s$, while norm takes values in $\operatorname{Cl}(L_i)[p]=0$.
Hence $s=-1$ on this module.
Every conjugate reflection also acts by $-1$, so their products,
which generate $C_3$, act trivially.
The nontrivial cubic class part is zero.
No assertion about the whole class group of $M_i$ is used.

\proofstep{4. The relevant global unit spaces have dimension two.}
The twelve complex places of $M_i$ above the chosen complex place
of $K$ form a regular $\Delta_i$-set.
The equivariant Dirichlet unit character is
$\mathbb Q[\Delta_i]-\mathbb Q$.
As $p\nmid12$ and $\mu_p(M_i)=1$, passage to units modulo $p$
preserves these character multiplicities.
For each $\lambda=82,187$, let $W_\lambda$ be the nontrivial
cubic part of the $b$-eigenspace with eigenvalue $\lambda$
in $\mathcal O_{M_i}^\times/\mathcal O_{M_i}^{\times p}$.
It is a two-dimensional $\mathbb F_p$-space.
Since $p\equiv2\bmod3$, the polynomial $X^2+X+1$ is
irreducible over $\mathbb F_p$, and $W_\lambda$ is a simple
$\mathbb F_p[C_3]$-module.
The localization map is $\Delta_i$-equivariant.
One nonzero local value from $W_\lambda$ therefore proves
injectivity on that whole space, and on its two cubic eigenlines
after extending scalars to $k$.

\proofstep{5. Two exact global units and their integral local conversion.}
Use unit number five for $L_1$ and unit number six for $L_2$
from the exact data file
\path{selected_units_exact.json}.
It specifies each element as
\[
         U_i=\frac{1}{18561}
           \sum_{0\leq a\leq2,\ 0\leq b\leq3}
                           n_{4a+b}\,t^a w_i^b .
 \tag{C16Q.23}
\]
The denominator is $18561=69p$.
The complete integer arrays and their monic integral
degree-twelve annihilators with constant coefficient one
are included in that file.  Its SHA256, split for printing, is
\begin{quote}\ttfamily\small
11993551e17b3919bc53a01a13fc1897\\
36775714559c19eacc47a57befa80844.
\end{quote}
The monic relations prove that $U_i$ are integral.
The reversed relations prove that their inverses are integral.
They are therefore actual global units.
No completeness of a candidate unit basis is required.

The independent rational verifier substitutes (C16Q.15) into
every tensor monomial, reduces by (C16Q.14), and only then
checks that every coefficient is $p$-integral.
It verifies the entire monic relation by Horner evaluation in
this rational degree-twelve ring.
It performs these checks on all fourteen supplied units, although
only the two displayed choices are used.
In particular it never discards an original denominator containing $p$.
The selected units reduce modulo $u^{12}$ to the following vectors
in the order $1,u,\ldots,u^{11}$:
\[
 \begin{aligned}
 \bar U_1&=(268,129,65,73,150,86,265,238,76,236,114,116),\\
 \bar U_2&=(268,64,57,156,57,164,142,6,249,182,67,260).
 \end{aligned}
 \tag{C16Q.24}
\]
Both constant terms are nonzero.
Normalize by the constant term and take the truncated logarithm.
All denominators $1,\ldots,11$ are units modulo $p$.
The two logarithm vectors are
\[
 \begin{aligned}
 L_1&=(0,140,88,127,0,231,12,233,60,227,192,211),\\
 L_2&=(0,205,47,185,0,81,177,56,78,238,127,171).
 \end{aligned}
 \tag{C16Q.25}
\]
The constant normalization causes no ambiguity in the next step,
because the cubic norm correction cancels constants.

\proofstep{6. Both quartic projectors, including the higher uniformizer terms.}
Form the global unit
$U_i^3/\prod_{r=0}^2d^r(U_i)$ before localization.
On its truncated logarithm, the operator
$3-(1+d+d^2)$ kills the coefficients with degree divisible by three
and multiplies all other coefficients by three.
For $\lambda=82,187$, apply the global operator
\[
                         P_\lambda=\sum_{r=0}^3\lambda^{-r}b^r
 \tag{C16Q.26}
\]
to this unit class modulo $p$.
The coefficients may be lifted to ordinary integer exponents.
It produces a class in $W_\lambda$.
Its local logarithm is computed using the actual action (C16Q.17);
in particular the $u^7$ term of $b(u)$ is retained.
The complete nonzero results modulo $u^{12}$ are
\[
 \begin{array}{c|c|l}
 \text{field and unit}&\lambda&\text{projected logarithm}\\ \hline
 L_1,\ U_1&82&218u^7+228u^{11}\\
 L_1,\ U_1&187&66u+82u^5+157u^7+152u^{11}\\
 L_2,\ U_2&82&39u+165u^5+164u^7+55u^{11}\\
 L_2,\ U_2&187&239u^7+114u^{11}.
 \end{array}
 \tag{C16Q.27}
\]
All entries lie in $\mathbb F_p$, embedded in the residue
field $\mathbb F_{p^2}$ of $M_i$.
The verifier checks both $b^2(u)=-u$ and
$b(P_\lambda L)=\lambda P_\lambda L$ at this precision.
The global operator and its localization agree because the
single $p$-place is stable under the entire $\Delta_i$.

Here the $p$-adic logarithm and exponential identify
$1+\mathfrak m_{M_p}$ with $\mathfrak m_{M_p}$.
Indeed $v_u(p)=12<p-1$, so beyond the linear term all logarithm
and exponential terms have larger valuation, and their formal
inverse identities converge.
Multiplication by $p$ sends $\mathfrak m_{M_p}$ to
$\mathfrak m_{M_p}^{13}$.
The residue roots of unity have order $p^2-1$, prime to $p$.
Thus the logarithm of a local $p$th power unit has no nonzero
term in the degrees displayed in (C16Q.27).
Every one of those four entries is therefore a genuine nonzero
local Kummer class.
The simple-module argument in Step4 proves injectivity on
both $W_\lambda$, and hence on every inverse-character unit
sector selected by (C16Q.19).
Together with the class-sector vanishing and (C16Q.22), this
proves the residual strict-zero assertion.

\proofstep{7. Lift through the original uniformizer.}
Let $\varpi$ be a uniformizer of the original $O$.
The module $D$ in (C16Q.5) is $\varpi$-divisible and
$D[\varpi]\simeq k(a)$ with the residual character already treated.
At $p$, $\phi$ is unramified, while the residual cubic inertia
has order three.
It cannot cancel any cyclotomic/Teichmuller residual factor,
whose order divides $268$; the reduction of $\kappa$ is trivial.
Hence some $g\in G_{K_p}$ has $\bar a(g)\ne1$, and
$D^{G_{K_p}}=0$.
At a finite place away from $p$, the cyclotomic, Teichmuller and
$\kappa$ factors are unramified.
The remaining inertia character has order dividing twelve,
prime to $p$.
Its invariants are either zero or the entire divisible module $D$.
Therefore in every required local group $J$ the connecting
quotient $D^J/\varpi D^J$ is zero, and
\[
                  H^1(J,D[\varpi])\longrightarrow H^1(J,D)
                         \quad\text{is injective}.
 \tag{C16Q.28}
\]
The global long exact sequence lifts every strict
$\varpi$-torsion class in $H^1(G_K,D)$ to $H^1(G_K,D[\varpi])$.
Local injectivity makes the lift strict, and the residual
strict kernel is zero.
Thus $\operatorname{Sel}_{\rm str}(K,D)[\varpi]=0$.
Every continuous cocycle has finite image in the discrete module
$E/O$, because $G_K$ is compact.
It is killed by a power of $\varpi$.
A nonzero strict class would therefore yield a nonzero
$\varpi$-torsion strict class, a contradiction.
This proves (C16Q.6) on the original ring, including ramified $O$.
\end{proof}

\paragraph{Actual finite labels and the precise weight count.}
For $m=0,\kappa=1$, the two new residues are
$e=134n+66$, $n\geq0$.
For a positive integer $T$, they contribute
\[
                 4\left\lfloor\frac{T+67}{134}\right\rfloor
 \tag{C16Q.29}
\]
triples $(\chi^{-j}\phi^{-\sigma},e)$ with $1\leq e<T$.
Combined with AC-G142 and AC-G144, all four residues are exactly
$e=67n-1$, $n\geq1$, and their count is $4\lfloor T/67\rfloor$.
The common-window/index comparison remains a separate application:
strict zero does not set the positive integral index to zero.
Every other finite label and coefficient prime is untouched here.

\paragraph{Source and geometric scope.}
AC-G147-class supplies the exact cyclotomic field identifications
and the prime-to-$269$ class conclusions, with its own complete
analytic and finite certificates.
The new local calculation uses the independent rational verifier
\path{verify_quartic_residue_local_units.py} and its full output
\path{quartic_residue_local_certificate.json}.
Kummer theory is the retained Stacks tag~03PK input;
prime-to-$p$ descent, global class extension/norm, and equivariant
Dirichlet units have the same source scope as AC-G130/134/142/144.
The finite-cover and rank-one local-system dictionary is the
retained SGA~1, Expos\'e~V, \S\S4 and 7 dictionary.
The extra quartic $p$-character labels a splitting cover for
Kummer descent; it does not change the original finite local system.
No new primary-source proof, complete fundamental-unit basis,
full class group of $M_i$, Lean execution, or mixed reciprocity
statement is claimed.
